W pracy przedstawiono przegląd metod kompresji wag modeli sztucznej inteligencji ze szczególnym uwzględnieniem technik kwantyzacji oraz zaprezentowano autorskie narzędzie "Uni-Quant" do szybkiej kwantyzacji i dekwantyzacji modeli TensorFlow.

Główne cele pracy obejmowały: analizę podstawowych pojęć związanych z reprezentacją liczb i kwantyzacją, omówienie współczesnych rozwiązań stosowanych w praktyce inżynierskiej oraz opracowanie praktycznej biblioteki umożliwiającej przeprowadzenie kwantyzacji i odzyskanie wag w formacie użytecznym do wnioskowania.

W części teoretycznej uporządkowano pojęcia tensora, parametrów modelu, formatów liczbowych oraz zagadnienia związane z kalibracją i metrykami oceny jakości modeli. Omówiono również wpływ wyboru reprezentacji na rozmiar modelu, przepustowość pamięciową i czas wykonywania obliczeń.

Rozdział dotyczący popularnych rozwiązań zawierał charakterystykę ekosystemu współczesnych narzędzi do kwantyzacji. Przedstawiono podejścia takie jak BitsAndBytes (w tym format NF4 i LLM.int8()), algorytmy optymalizujące błąd rekonstrukcji wyjść (GPTQ), podejścia aktywacyjne uwzględniające ważność kanałów (AWQ) oraz nowsze metody typu TurboQuant. Omówiono także techniki redukcji metadanych, takie jak podwójna kwantyzacja, oraz formaty plików i praktyczne ograniczenia wdrożeń produkcyjnych.

W części implementacyjnej zaprezentowano bibliotekę Uni-Quant, która realizuje post-training quantization w wariantach 4- i 8-bitowych, opierając się na grupowaniu wag w bloki i zapisie współczynnika skalującego dla każdego bloku. Wprowadzono mechanizmy:
\begin{itemize}
    \item Blokowego skalowania wag oraz zapisu skal,\\
    \item Pakowania dwóch wartości 4-bitowych w jeden bajt,\\
    \item Zapisu wag i metadanych w postaci binarnej oraz JSON-owych manifestów scalonych w archiwum z rozszerzeniem .uniq,\\
    \item Zachowania fragmentów o rozmiarze mniejszym niż blok w postaci float32, aby zminimalizować degradację jakości.
\end{itemize}

Dodatkowo opracowano alternatywną implementację z wykorzystaniem CUDA, przenosząc krytyczne obliczenia do kerneli uruchamianych na GPU w celu przyspieszenia procesu kwantyzacji i dekwantyzacji. Rozwiązanie to umożliwia znaczące skrócenie czasu przetwarzania dużych modeli i stanowi praktyczny wybór tam, gdzie dostępne są akceleratory GPU.

Testy przeprowadzone w pracy wykazały oczekiwane zależności: zmniejszenie liczby bitów prowadzi do istotnego spadku rozmiaru pliku modelu, natomiast wpływ na jakość zależy od strategii grupowania wag (rozmiar bloku) oraz szczegółów kalibracji. Mniejszy rozmiar bloku zwykle zwiększa liczbę współczynników skalowania, co może prowadzić do większego rozmiaru pliku przy bardzo drobnej granularności, ale pozwala lepiej dopasować skalę lokalną i ograniczyć stratę jakości. Zachowanie niektórych małych fragmentów wag w formacie FP32 okazało się praktycznym kompromisem redukującym negatywny wpływ kwantyzacji na krytyczne parametry modelu.

Własny wkład praktyczny można podsumować jako:
\begin{itemize}
    \item Zaprojektowanie i zaimplementowanie biblioteki Uni-Quant do szybkiej kwantyzacji i dekwantyzacji modeli TensorFlow wraz z formatem archiwum .uniq,\\
    \item Przeprowadzenie testów ilustrujących kompromisy między rozmiarem pliku a zachowaniem jakości oraz wpływu rozmiaru bloku i szerokości bitowej,\\
    \item Integracja przyspieszenia obliczeń opartego na kernelach CUDA oraz udokumentowanie wpływu akceleracji na czasy przetwarzania.
\end{itemize}

Ograniczenia pracy i kierunki dalszych badań: praca koncentruje się na podejściach post-training quantization i prostych strategiach blokowego skalowania. Do istotnych rozszerzeń należą integracja zaawansowanych algorytmów PTQ (np. GPTQ i AWQ), implementacja per-channel scaling i technik kalibracji opartych na reprezentatywnych zbiorach danych. Dalsze badania mogą również objąć automatyczne dobieranie parametrów (np. rozmiar bloku) przy uwzględnieniu kosztu pamięciowego i metryk jakości oraz bardziej rozbudowane testy na modelach produkcyjnych o bardzo dużej liczbie parametrów. Praktycznym kierunkiem rozwoju jest bezpośrednia integracja formatu .uniq z ekosystemem TensorFlow, umożliwiająca wykonywanie modelu bezpośrednio z archiwum .uniq przy użyciu mechanizmu on‑the‑fly (lub częściowej, leniwej) dekwantyzacji bloków podczas inferencji, co poza obecnym zmniejszaniem rozmiaru na dysku, obniży wymagania sprzętowe i pozwoli uruchamiać duże modele w środowiskach o ograniczonych zasobach.

Podsumowując, praca ukazuje, że kwantyzacja jest efektywnym mechanizmem zmniejszania kosztów przechowywania i wykonywania modeli SI, lecz wymaga starannego doboru strategii i parametrów, aby zminimalizować spadek jakości. Uni-Quant dostarcza praktycznego narzędzia do eksperymentów i stanowi solidną bazę do dalszego rozwoju metod kompresji.
